\subsection{UV seed-shadow identities and the Cabibbo inversion}

The retained scalar seed $\phiz$ still satisfies exact UV identities. In the present version these
identities are auxiliary kernel bookkeeping rather than the independence grammar of the physical
observable layer. In particular, the Cabibbo angle admits an explicit inversion back to the UV
seed:
\begin{equation}
\label{eq:cabibbo-inversion}
\phiz=\frac{1-\sqrt{1-4\lambda_C^2}}{2}.
\end{equation}
\begin{remark}[Catalan expansion of the Cabibbo inversion]
Setting
\[
x:=\lambda_C^2,
\]
the inversion formula becomes
\[
\phiz=x\,C(x),
\qquad
C(x):=\frac{1-\sqrt{1-4x}}{2x}.
\]
Hence
\begin{equation}
\label{eq:catalan-cabibbo}
\phiz
=
\sum_{n\ge 0} C_n\,\lambda_C^{2n+2}
=
\lambda_C^2+\lambda_C^4+2\lambda_C^6+5\lambda_C^8+14\lambda_C^{10}+\cdots,
\end{equation}
where $C_n$ are the Catalan numbers. In the present paper the appropriate interpretation of this
identity is noncrossing combinatorics on the seed side; it is a theorem-level strengthening of
the one-seed compression rather than a claim about planar Feynman graphs, a specific large-$N$ /
't Hooft resummation, or a holographic flavor limit.
\end{remark}

\begin{proposition}[Retained seed as one decoder]
\label{prop:retained-seed-decoder}
Set
\[
u:=\phiz.
\]
On the retained branch one has
\[
\lambda_C^2=u(1-u),
\qquad
\beta_{\mathrm{rad}}=\frac{u}{4\pi},
\qquad
\Omega_b=\left(1-\frac{1}{4\pi}\right)u,
\qquad
\sin^2\theta_{13}=e^{-\gamma}u.
\]
Equivalently,
\[
u
=
\frac{1-\sqrt{1-4\lambda_C^2}}{2}
=
4\pi\beta_{\mathrm{rad}}
=
\frac{\Omega_b}{1-\frac{1}{4\pi}}
=
e^\gamma\sin^2\theta_{13}.
\]
Thus the retained branch carries a single decoder
\[
u\longmapsto (\lambda_C,\beta_{\mathrm{rad}},\Omega_b,\theta_{13}).
\]
\end{proposition}

\begin{proof}
The forward identities are exactly the retained-seed bridge relations. The inverse identities follow
from the Cabibbo inversion \eqref{eq:cabibbo-inversion} and direct elimination of $u$ from the same
four formulas.
\end{proof}

Eliminating the retained seed through the Cabibbo inversion gives the equivalent
$\lambda_C$-only formulas:
\begin{equation}
\label{eq:lambda-c-master}
\begin{aligned}
\beta_{\mathrm{rad}}&=\frac{1-\sqrt{1-4\lambda_C^2}}{8\pi},\\
\Omega_b&=\frac{4\pi-1}{8\pi}\bigl(1-\sqrt{1-4\lambda_C^2}\bigr),\\
\sin^2\theta_{13}&=\frac{e^{-\gamma}}{2}\bigl(1-\sqrt{1-4\lambda_C^2}\bigr).
\end{aligned}
\end{equation}

\begin{corollary}[Vacuum half-circle and UV-angle parametrization]
\label{cor:uv-half-circle}
On the retained UV seed split one has the exact identity
\[
\lambda_C^2+\left(\phiz-\frac12\right)^2=\frac14.
\]
Defining
\[
\sin^2\theta_{\mathrm{UV}}:=\phiz,
\]
one obtains
\[
\lambda_C=\frac12\sin(2\theta_{\mathrm{UV}}),
\qquad
\sin^2\theta_{13}=e^{-\gamma}\sin^2\theta_{\mathrm{UV}},
\qquad
\Omega_b=\left(1-\frac{1}{4\pi}\right)\sin^2\theta_{\mathrm{UV}}.
\]
\end{corollary}

\begin{proof}
The exact seed relation $\lambda_C^2=\phiz(1-\phiz)$ is equivalent to the displayed half-circle
identity after completing the square. The $\theta_{\mathrm{UV}}$ parametrization then follows
immediately from $\sin^2\theta_{\mathrm{UV}}=\phiz$ together with the exact retained-seed formulas
for $\sin^2\theta_{13}$ and $\Omega_b$.
\end{proof}

\begin{corollary}[CKM--PMNS bridge law]
\label{cor:ckm-pmns-bridge-law}
On the closed branch one has the exact relation
\[
\sin^2\theta_{13}
=
\frac{e^{-5/6}}{2}\Bigl(1-\sqrt{1-4\lambda_C^2}\Bigr).
\]
Equivalently,
\[
\lambda_C^2=e^{5/6}\sin^2\theta_{13}-e^{10/6}\sin^4\theta_{13}.
\]
In internal branch form,
\[
\sin^2\theta_{13}=0.0231084351588811\ldots
\]
with no additional fit parameter.
\end{corollary}

\begin{proof}
By \Cref{eq:cabibbo-inversion},
\[
\phiz=\frac{1-\sqrt{1-4\lambda_C^2}}{2}.
\]
By the exact seed identity,
\[
\sin^2\theta_{13}=\phiz e^{-\gamma},
\qquad
\gamma=\frac56.
\]
Substituting $\phiz$ gives the displayed formula.
\end{proof}

\begin{corollary}[Baryon--reactor bridge law]
\label{cor:baryon-reactor-bridge-law}
On the closed branch one has the exact cross-sector relation
\[
\Omega_b
=
\left(1-\frac{1}{4\pi}\right)e^{5/6}\sin^2\theta_{13}
=
\frac{4\pi-1}{8\pi}\Bigl(1-\sqrt{1-4\lambda_C^2}\Bigr).
\]
In internal branch form,
\[
\Omega_b=0.0489406626654501\ldots
\]
with no additional fit parameter.
\end{corollary}

\begin{proof}
The exact seed-budget identities give
\[
\Omega_b=(4\pi-1)\beta_{\mathrm{rad}},
\qquad
\beta_{\mathrm{rad}}=\frac{\phiz}{4\pi},
\qquad
\sin^2\theta_{13}=\phiz e^{-5/6}.
\]
Eliminating $\phiz$ yields the displayed cross-sector identity. Eliminating $\phiz$ instead through
\Cref{eq:cabibbo-inversion} gives the equivalent $\lambda_C$ form.
\end{proof}

\begin{corollary}[Matter-radiation quotient on the retained seed split]
\label{cor:matter-radiation-quotient}
On the retained seed split one has the exact quotient identity
\[
\frac{\Omega_b}{\beta_{\mathrm{rad}}}=4\pi-1.
\]
\end{corollary}

\begin{proof}
This is the direct quotient of the exact retained-seed identities
\[
\Omega_b=(4\pi-1)\beta_{\mathrm{rad}},
\qquad
\beta_{\mathrm{rad}}=\frac{\phiz}{4\pi}.
\]
\end{proof}

\begin{remark}[Interpretive reading of the retained seed split]
The same one-seed decoder carries the exact budget identity
\begin{equation}
\label{eq:seed-budget-identity}
\Omega_b+\beta_{\mathrm{rad}}=\phiz.
\end{equation}
Indeed,
\[
\Omega_b+\beta_{\mathrm{rad}}
=(4\pi-1)\beta_{\mathrm{rad}}+\beta_{\mathrm{rad}}
=4\pi\,\beta_{\mathrm{rad}}
=\phiz.
\]
Both \Cref{cor:matter-radiation-quotient,eq:seed-budget-identity} are exact retained-seed
identities. Any solid-angle or spherical-surface interpretation remains interpretive and is not
used as theorem input.
\end{remark}

Thus on the canonical branch, $\beta$, $\Omega_b$, $\lambda_C$, and $\theta_{13}$ are four exact
projections of the same retained decoder $u=\phiz$. Conversely, $\theta_{13}$ fixes $\lambda_C$ and
$\Omega_b$ through $\phiz=e^\gamma\sin^2\theta_{13}$. This is the derived-consequence layer of the
retained seed split, not a separate independent theorem surface.

\subsection{Carrier comparison of alternative discrete worlds}

The strongest numerical pressure test on the carrier packet is not a further precision hit. It is
that nearby wrong discrete worlds miss the electromagnetic closure decisively.

\begin{proposition}[Numerical comparison of alternative discrete worlds]
\label{thm:carrier-exclusion}
Let
\[
\alpha^{-1}(0;N_{\mathrm{fam}},N_\Phi,\nu^c)
\]
denote the positive self-consistent root of the carrier-form electromagnetic closure equation
\eqref{eq:carrier-cfe} for the indicated discrete data. Then:
\begin{center}
\small
\sisetup{round-mode=none, group-minimum-digits=20}
\setlength{\tabcolsep}{2pt}
\renewcommand{\arraystretch}{0.96}
\begin{tabularx}{0.98\linewidth}{@{}>{\raggedright\arraybackslash}p{0.21\linewidth}>{\raggedright\arraybackslash}p{0.19\linewidth}>{\raggedleft\arraybackslash}p{0.14\linewidth}>{\raggedleft\arraybackslash}p{0.14\linewidth}Y@{}}
\toprule
Discrete world & Change from canonical & $\alpha^{-1}(0)$ & Residual vs CODATA & Reading \\
\midrule
\textbf{canonical $(3,1,\text{yes})$} & --- & $\mathbf{\num{137.0359992}}$ & $+\num{3.9e-8}$ & \makecell[l]{inside CODATA window\\($+1.87\sigma$)} \\
$N_{\mathrm{fam}}=2$ & one family fewer & $\num{156.094}$ & $+19.1$ & catastrophically excluded \\
$N_{\mathrm{fam}}=4$ & one family more & $\num{124.835}$ & $-12.2$ & catastrophically excluded \\
$N_\Phi=2$ & two light Higgs doublets & $\num{135.946}$ & $-1.09$ & catastrophically excluded \\
no $\nu^c$ & remove right-handed neutrino & $\num{137.0338}$ & $-\num{2.16e-3}$ & \makecell[l]{excluded\\(${\sim}\num{1.03e5}\sigma$)} \\
\bottomrule
\end{tabularx}
\end{center}
\end{proposition}

\begin{proof}[Proof sketch]
Hold fixed
\[
\cthree=\frac{1}{8\pi},
\qquad
\phibase=\frac{1}{6\pi},
\qquad
\gamma=\frac56,
\qquad
B=\frac32,
\]
and solve the positive self-consistent root of the same scalar closure equation on each discrete world:
\[
\alpha^3-2\cthree^3\alpha^2-8\cthree^6\,\bone(N_{\mathrm{fam}},N_\Phi)\ln\!\bigl(\phiseam(\alpha)^{-1}\bigr)=0,
\]
with
\[
\phiseam(\alpha)=\phibase+\deltatop e^{-2\alpha}(1-\deltatop e^{-2\alpha})^{-5/4},
\qquad
\deltatop=\Oadm\cthree^4,
\qquad
\delta_2=\frac54\,\deltatop^2,
\]
and
\[
\bone(N_{\mathrm{fam}},N_\Phi)=\frac43N_{\mathrm{fam}}+\frac1{10}N_\Phi.
\]
For the no-$\nu^c$ row, $\bone$ stays unchanged because $Y(\nu^c)=0$, while the occupancy changes from $\Oadm=48$ to $\Oadm=45$ because the one-family packet drops from $16$ to $15$ states. The same positive-branch solver and tolerance are used for every row, so the deviations are arithmetic consequences of changed discrete data rather than fits.
\end{proof}

\begin{remark}
Two families shift $\alpha^{-1}$ by $+19$, four families by $-12$, and two Higgs doublets by
$-1.09$. Even the mildest alternative---removing $\nu^c$---shifts the value by
$\num{2.16e-3}$, which is $\sim\num{1.03e5}$ times the CODATA uncertainty. This is therefore a
useful comparison pressure test on the canonical branch rather than an additional theorem in the
reconstruction chain.
\end{remark}

\begin{figure}[t]
\centering
\scriptsize
\begin{tikzpicture}[
    >=Latex,
    head/.style={draw=black!65, fill=gray!15, rounded corners, minimum height=0.8cm, align=center, font=\bfseries},
    ok/.style={draw=green!50!black, fill=green!8, rounded corners, minimum height=0.8cm, align=left},
    bad/.style={draw=red!70!black, fill=red!6, rounded corners, minimum height=0.8cm, align=left}
]
\node[head, minimum width=3.5cm] at (0,0) {discrete world};
\node[head, minimum width=2.8cm] at (4.15,0) {Residual vs CODATA};
\node[head, minimum width=4.6cm] at (8.2,0) {reading};

\node[ok, anchor=west, text width=3.25cm] at (-1.72,-1.0) {canonical $(3,1,\mathrm{yes})$};
\node[ok, minimum width=2.55cm] at (4.15,-1.0) {$+\num{3.9e-8}$};
\node[ok, anchor=west, text width=4.35cm] at (6.0,-1.0) {inside the CODATA window};

\node[bad, anchor=west, text width=3.25cm] at (-1.72,-1.95) {$N_{\mathrm{fam}}=2$};
\node[bad, minimum width=2.55cm] at (4.15,-1.95) {$+19.1$};
\node[bad, anchor=west, text width=4.35cm] at (6.0,-1.95) {catastrophically excluded};

\node[bad, anchor=west, text width=3.25cm] at (-1.72,-2.9) {$N_{\mathrm{fam}}=4$};
\node[bad, minimum width=2.55cm] at (4.15,-2.9) {$-12.2$};
\node[bad, anchor=west, text width=4.35cm] at (6.0,-2.9) {catastrophically excluded};

\node[bad, anchor=west, text width=3.25cm] at (-1.72,-3.85) {$N_\Phi=2$};
\node[bad, minimum width=2.55cm] at (4.15,-3.85) {$-1.09$};
\node[bad, anchor=west, text width=4.35cm] at (6.0,-3.85) {extra light doublet breaks closure};

\node[bad, anchor=west, text width=3.25cm] at (-1.72,-4.8) {no $\nu^c$};
\node[bad, minimum width=2.55cm] at (4.15,-4.8) {$-\num{2.16e-3}$};
\node[bad, anchor=west, text width=4.35cm] at (6.0,-4.8) {$\sim\num{1.03e5}\sigma$ exclusion};
\end{tikzpicture}
\caption{Comparison matrix for alternative discrete worlds. The middle column records the residual
versus the CODATA benchmark rather than a misleading shift relative to the canonical row. Once the
discrete packet is changed, the electromagnetic closure leaves the observed window decisively.}
\label{fig:carrier-exclusion-matrix}
\end{figure}

\section{Closure architecture A: families, Higgs, and local spectral scale}

This first closure-architecture block is the shortest route from the hard carrier packet to the
closed theorem body. It has three moves: primitive winding balance in the seam normal plane forces
the four-corner structure of the fundamental section, that winding theorem globalizes into the
family space $F$, the determinant line selects the unique light Higgs doublet together with the
compact seam phase, and the local spectral-scale channel generates both $\Mpl$ and the electroweak
vacuum from one common geometric response.

